% !TITLE 诗经·卫风·淇奥
% !AUTHOR 无名氏
% !DYNASTY 先秦
% !GENRE 四言诗
% !ORDER 3
% !SOURCE 诗经·卫风

\begin{poem}
瞻彼淇奥，绿竹猗猗。有匪君子，如切如磋，如琢如磨。
瑟兮僩兮，赫兮咺兮。有匪君子，终不可谖兮。

瞻彼淇奥，绿竹青青。有匪君子，充耳琇莹，会弁如星。
瑟兮僩兮，赫兮咺兮。有匪君子，终不可谖兮。

瞻彼淇奥，绿竹如箦。有匪君子，如金如锡，如圭如璧。
宽兮绰兮，猗重较兮。善戏谑兮，不为虐兮。
\end{poem}

\begin{appreciation}
如果说《关雎》写的是爱情的开端，《鹿鸣》写的是友朋的宴饮，那么《淇澳》写的便是一个理想的人。

淇水之岸，绿竹猗猗。竹子挺拔、有节、中空、长青，这四个特质恰恰是中国文化里对君子的全部想象：挺立而不屈，有节操而不放纵，虚怀若谷，始终不改其色。诗开篇以竹起兴，之后便直接描绘君子："如切如磋，如琢如磨"——他像一块玉石，经过反复的切割和打磨才达到如此温润的境界。这话里有一个重要的前提：没有人天生就是君子。修身是一个漫长的、需要忍耐的过程，就如治玉，一刀一刀，磨去粗粝，见出光泽。

后面又说"如金如锡，如圭如璧"。金与锡是可熔铸的，圭与璧是玉制的礼器——君子既要有可塑的质地，也要有不可轻慢的威仪。但最妙的是末章那两句："宽兮绰兮""善戏谑兮，不为虐兮"。他宽厚温和，甚至还会开玩笑——只是不过分。一个整天板着脸讲道德的人，是无趣的。真正的君子，温润如玉，近之可亲。玉的美，不是夺目耀眼的那种，而是温润内敛，摸上去有温度，看起来有光泽。竹子的美，也不是花朵那种招摇的，而是清俊挺拔，风吹过会响。这两种事物在《淇澳》里交织成同一个比兴，用来写君子，再恰当不过。

《诗经》三百篇，写人的诗很多，但像《淇澳》这样完整勾勒出一个理想人格轮廓的，不多。后世文人张口就来"谦谦君子，温润如玉"，源头便在此处。
\end{appreciation}
